% Part: lambda-calculus
% Chapter: introduction
% Section: syntax

\documentclass[../../../include/open-logic-section]{subfiles}

\begin{document}

\olfileid{lam}{int}{syn}
\olsection{ল্যাম্বডা ক্যালকুলাসের বাক্যগঠন}

প্রথমে $x$, $y$, $z$,~\dots চলরাশির একটি অনুক্রম এবং $a$, $b$,
$c$,~\dots কয়েকটি ধ্রুবক প্রতীক নেওয়া হয়। পদের সমষ্টি নিচের মতো
আরোহীভাবে সংজ্ঞায়িত:
\begin{enumerate}
\item প্রতিটি চলরাশি একটি পদ।
\item প্রতিটি ধ্রুবক একটি পদ।
\item $M$ ও $N$ পদ হলে $(MN)$-ও একটি পদ।
\item $M$ একটি পদ এবং $x$ একটি চলরাশি হলে $(\lambd[x][M])$ একটি পদ।
\end{enumerate}
$(MN)$ রূপের পদগুলিকে \emph{প্রয়োগ-পদ} এবং $(\lambd[x][M])$ রূপের
পদগুলিকে \emph{বিমূর্তন-পদ} বলে।

কোনো ধ্রুবকই নেই এমন ব্যবস্থাকে \emph{বিশুদ্ধ} ল্যাম্বডা ক্যালকুলাস
বলে। আমরা প্রধানত বিশুদ্ধ $\lambd$-ক্যালকুলাসে কাজ করব, তাই সব ছোট হাতের
অক্ষর চলরাশি নির্দেশ করবে। বড় হাতের অক্ষর ($M$, $N$ ইত্যাদি) দিয়ে আমরা
$\lambd$-ক্যালকুলাসের পদ নির্দেশ করি।

আমরা কয়েকটি সাংকেতিক প্রথা অনুসরণ করব:

\begin{conv}
\begin{enumerate}
\item প্রথম বন্ধনী বাদ গেলে প্রয়োগ বাম থেকে ডানে ঘটে। যেমন, $M$, $N$,
  $P$ ও $Q$ পদ হলে $MNPQ$ হলো $(((MN)P)Q)$-এর সংক্ষিপ্ত রূপ।
\item আবার, প্রথম বন্ধনী বাদ গেলে ল্যাম্বডা-বিমূর্তনকে সম্ভাব্য সর্বাধিক
  পরিসর দিতে হবে। যেমন, $\lambd[x][MNP]$-কে
  $(\lambd[x][((MN)P)])$ পড়া হয়।
\item একটি ল্যাম্বডা দিয়ে একাধিক চলরাশিকে বিমূর্ত করা যায়। যেমন,
  $\lambd[xyz][M]$ হলো
  $\lambd[x][\lambd[y][\lambd[z][M]]]$-এর সংক্ষিপ্ত রূপ।
\end{enumerate}
\end{conv}

যেমন,
\[
\lambd[xy][xxyx \lambd[z][xz]]
\]
হলো নিচেরটির সংক্ষিপ্ত রূপ:
\[
\lambd[x][\lambd[y][((((xx)y)x)(\lambd[z][(xz)]))]].
\]
এই প্রথাগুলি মুখস্থ করে নাও। প্রথমে এগুলি তোমাকে
নাজেহাল করবে, কিন্তু ক্রমে অভ্যস্ত হয়ে যাবে; আর কিছুদিন পরে প্রথম বন্ধনীর
জঙ্গল সামলানোর চেয়ে এগুলিই কম নাজেহাল করবে।

শুধু বদ্ধ চলরাশির নামে তফাত আছে এমন দুটি পদের সম্পর্ককে
\emph{আলফা-সমতুল্যতা} বলে; অর্থাৎ পদদুটি $\alpha$-সমতুল্য। যেমন,
$\lambd[x][x]$ ও $\lambd[y][y]$। এগুলিকে ``একই'' পদ মনে করা সুবিধাজনক;
অন্যভাবে বললে, $M$ ও $N$ একই বললে আমরা ``বদ্ধ চলরাশিগুলির নামবদল পর্যন্ত
একই''-ও বোঝাব। কোনো $\lambd$-র পরিসরের মধ্যে থাকা চলরাশিগুলিকে ``বদ্ধ''
এবং অন্যগুলিকে ``মুক্ত'' বলে। আগের উদাহরণে কোনো মুক্ত চলরাশি নেই; কিন্তু
\[
(\lambd[z][yz])x
\]
-এ $y$ ও $x$ মুক্ত এবং $z$ বদ্ধ।

\end{document}
